% =====================================================================
%  BÖLÜM 3: Rasgele Değişkenler ve Dağılımlar
%  Kaynak: Feller Vol.I-II; Shiryaev; Billingsley; Klenke
% =====================================================================
\chapter{Rasgele Değişkenler ve Dağılımlar}
\label{chp:rasgele_degiskenler}

\bolumalinti{%
  Olasılık teorisi, belirsizliği sayıya dönüştürmenin bilimidir.
  Rasgele değişken bu dönüşümün ta kendisidir.%
}{Andrei Kolmogorov}

% ---- Kazanımlar (TYMM) ----
\begin{kazanimlar}
Bu bölümün sonunda öğrenci:
\begin{itemize}[leftmargin=1.8em]
  \item Rasgele değişkeni ölçülebilir fonksiyon olarak tanımlayabilir ve
        bu tanımın neden "şans sonucu sayı üretme" sezgisini yakaladığını açıklayabilir.
  \item Kümülatif dağılım fonksiyonu (KDF) kavramını tanımlayabilir,
        karakteristik özelliklerini ispatlayabilir ve KDF'den olasılık hesabı yapabilir.
  \item Kesikli, mutlak sürekli ve karma dağılımlar arasındaki farkı
        Lebesgue ayrışımı çerçevesinde ifade edebilir.
  \item Ortak dağılım, marjinal dağılım ve koşullu dağılım kavramlarını
        tanımlayabilir; bağımsız rasgele değişkenlerin ayırt edici özelliklerini kullanabilir.
  \item Birebir dönüşüm yöntemi ve Jacobian formülünü tek ve çok boyutlu
        durumlarda uygulayabilir; evrişim (konvolüsyon) yöntemiyle toplam dağılımını bulabilir.
\end{itemize}
\end{kazanimlar}

% ---- Motivasyon ----
\begin{motivasyon}
Bölüm~\ref{chp:olasilik_uzaylari}'da bir olasılık uzayı $(\Omega, \mathcal{F}, P)$ inşa ettik.
Ancak $\Omega$ kümesi çoğu zaman soyut bir uzaydır: yazı-tura atışında
$\Omega = \{Y, T\}$, bir deneyin tüm sonuçları kümesinde ise karmaşık nesneler barındırabilir.
Pratikte ilgilendiğimiz şey genellikle \emph{sayısal} sonuçlardır:
"kaç tura geldi?", "bekleme süresi ne kadar?", "ölçülen sıcaklık nedir?"

\textbf{Rasgele değişken} (random variable), $\Omega$'yı reel sayı eksenine taşıyan
ölçülebilir bir fonksiyondur. Bu taşıma sayesinde karmaşık olasılık uzaylarını
\emph{gerçel sayılar üzerindeki dağılımlarla} inceleyebiliriz; asıl ağırlık
$(\Omega, \mathcal{F}, P)$'den $(\mathbb{R}, \mathcal{B}(\mathbb{R}), P_X)$'e geçer.

\medskip
\textbf{Köprü.} Bölüm~\ref{chp:olcum_teorisi}'de ölçülebilir fonksiyon ve
Lebesgue integrali; Bölüm~\ref{chp:olasilik_uzaylari}'da Kolmogorov aksiyomları
ve bağımsızlık kavramı bu bölümün temelini oluşturur.
Bölüm~\ref{chp:beklenti}'de bu dağılımlar üzerinden beklenti ve momentler hesaplanacak;
Bölüm~\ref{chp:dagilim_aileleri}'de somut dağılım aileleri incelenecektir.
\end{motivasyon}

% =====================================================================
\section{Rasgele Değişken Kavramı}
\label{sec:rd_kavram}
% =====================================================================

\begin{kesif}
Yazı-tura atışında $\Omega = \{Y, T\}$ olsun.
"Tura sayısı" bir fonksiyon tanımlar: $X(Y) = 0$, $X(T) = 1$.
Bu fonksiyonun hangi özelliği onu "rasgele değişken" yapar?
Eğer $A = \{1\} \subset \mathbb{R}$ ise $\{X \in A\} = \{T\}$ olayı
olasılık uzayında tanımlı olmalı — yani $\mathcal{F}$'ye ait olmalıdır.
Peki ya $A = (0.5, 1.5)$ veya $A = \emptyset$ için?
Tüm Borel kümelerine bu özelliği istiyorsak ne elde ederiz?
\end{kesif}

\begin{tanimgerekce}
Neden "$f^{-1}(B) \in \mathcal{F}$ her $B \in \mathcal{B}(\mathbb{R})$ için" koşulu?
Çünkü $P(X \in B)$ ifadesi ancak $\{X \in B\} = X^{-1}(B)$ bir olay olduğunda,
yani $\mathcal{F}$'de olduğunda, anlam taşır.
Ölçülebilirlik tam da bu hesabı mümkün kılan koşuldur.
\end{tanimgerekce}

\begin{tanim}[Rasgele Değişken]
\label{def:rasgele_degisken}
$(\Omega, \mathcal{F}, P)$ bir olasılık uzayı olsun.
$X : \Omega \to \mathbb{R}$ fonksiyonu, her $B \in \mathcal{B}(\mathbb{R})$ için
\[
  X^{-1}(B) = \{\omega \in \Omega : X(\omega) \in B\} \in \mathcal{F}
\]
koşulunu sağlıyorsa $X$'e \textbf{rasgele değişken} (random variable) denir.
Başka bir deyişle $X$, $(\Omega, \mathcal{F})$'den $(\mathbb{R}, \mathcal{B}(\mathbb{R}))$'ye
\textbf{ölçülebilir bir fonksiyondur}.
\end{tanim}

\begin{onerme}[Eşdeğer Karakterizasyon]
\label{prop:rd_char}
$X : \Omega \to \mathbb{R}$ için aşağıdaki koşullar birbirine denktir:
\begin{enumerate}[(i)]
  \item $X$ rasgele değişkendir.
  \item Her $c \in \mathbb{R}$ için $\{X \leq c\} \in \mathcal{F}$.
  \item Her $c \in \mathbb{R}$ için $\{X < c\} \in \mathcal{F}$.
  \item Her $c \in \mathbb{R}$ için $\{X > c\} \in \mathcal{F}$.
  \item Her $c \in \mathbb{R}$ için $\{X \geq c\} \in \mathcal{F}$.
\end{enumerate}
\end{onerme}

\begin{proof}
Bölüm~\ref{chp:olcum_teorisi}'deki Teorem~\ref{thm:olculebichar}'in doğrudan
bir uygulamasıdır; $\pi$-sistemi argümanı $\{(-\infty, c] : c \in \mathbb{R}\}$'nin
$\mathcal{B}(\mathbb{R})$'yi ürettiğini kullanır.
$(i) \Rightarrow (ii)$: $(-\infty, c] \in \mathcal{B}(\mathbb{R})$ olduğundan $X^{-1}((-\infty, c]) \in \mathcal{F}$.
$(ii) \Rightarrow (iii)$: $\{X < c\} = \bigcup_{n=1}^\infty \{X \leq c - 1/n\}$,
sayılabilir birleşim $\mathcal{F}$'de kalır.
$(iii) \Leftrightarrow (iv)$: $\{X < c\} = \{X \geq c\}^c$.
$(ii) \Leftrightarrow (v)$: $\{X \geq c\} = \{X < c\}^c$.
$(ii) \Rightarrow (i)$: $\pi$-sistemi $\{(-\infty, c]\}$ üzerinde
$X^{-1}$ $\mathcal{F}$'ye düşüyor; $\mathcal{B}(\mathbb{R})$'nin üretilmesi bu yeterli.
\end{proof}

\begin{kritiksorgu}
"$X$ sürekli olması rasgele değişken olması için yeterli midir?"
Hayır: süreklilik analitik bir kavramdır, ölçülebilirliği garanti etmez —
ancak Borel-ölçülebilir bir $\sigma$-cebir düzeni altında sürekli fonksiyonlar
her zaman ölçülebilirdir (dolayısıyla rasgele değişkendir).
Vitali kümesinin gösterge fonksiyonu, $(\mathbb{R}, \mathcal{B}(\mathbb{R}))$'de
ölçülebilir olmayan bir fonksiyon örneğidir (ama bu durum uygulamada sorun yaratmaz).
\end{kritiksorgu}

\begin{onerme}[Cebirsel Kararlılık]
\label{prop:rd_cebir}
$X, Y$ rasgele değişkenler ve $c \in \mathbb{R}$ sabit olsun. O zaman
$X + Y$, $cX$, $XY$, $\max(X, Y)$, $\min(X, Y)$ ve $|X|$ de rasgele değişkenlerdir.
Ayrıca $Y \neq 0$ n.k.\ ise $X/Y$ rasgele değişkendir.
$\{X_n\}$ rasgele değişken dizisi ise $\sup_n X_n$, $\inf_n X_n$,
$\limsup_{n} X_n$, $\liminf_n X_n$ ölçülebilir fonksiyonlardır
(sonlu değerleri yoksa $\overline{\mathbb{R}}$'ye değer alırlar).
\end{onerme}

\begin{proof}
$\{X + Y < c\} = \bigcup_{q \in \mathbb{Q}} (\{X < q\} \cap \{Y < c-q\}) \in \mathcal{F}$
(rasyonel sayılar üzerinden sayılabilir birleşim).
$\{XY > c\}$ için $XY = \tfrac{1}{4}[(X+Y)^2 - (X-Y)^2]$ açılımı kullanılır.
$\sup_n$: $\{\sup_n X_n > c\} = \bigcup_n \{X_n > c\} \in \mathcal{F}$.
Diğerleri benzer.
\end{proof}

\begin{tarihselnot}
"Rasgele değişken" terimi bir değişken değil, fonksiyondur — tarihsel adlandırma yanıltıcıdır.
Kolmogorov 1933'te olasılığı ölçüm teorisi üzerine oturtmadan önce,
rasgele değişkenler daha çok sezgisel olarak tanımlanıyordu.
Modern tanım, Lebesgue'nin ölçüm teorisi ile Borel'in $\sigma$-cebirleri
kavramlarının olasılığa uygulanmasıyla ortaya çıktı.
\end{tarihselnot}

\begin{hocanotu}
Bundan sonra "$X$ rasgele değişken" derken her zaman
$(\Omega, \mathcal{F}, P)$ olasılık uzayından $(\mathbb{R}, \mathcal{B}(\mathbb{R}))$'ye
ölçülebilir fonksiyonu kastediyoruz. Vektör değerli rasgele değişkenler
(Bölüm~\ref{sec:rd_ortak}) aynı mantıkla $\mathbb{R}^n$'e değer alan ölçülebilir
fonksiyonlar olarak tanımlanır.
\end{hocanotu}

\begin{kendinleifade}
"Rasgele değişken" neden bir "değişken" değil, bir "fonksiyondur"?
Hangi uzaydan hangi uzaya gider? Ölçülebilirlik koşulu ne anlama gelir?
Bir arkadaşınıza matematiksiz bir dille açıklayın.
\ogrencialani[2.5cm]
\end{kendinleifade}

\begin{mikroozet}
Rasgele değişken $X : (\Omega, \mathcal{F}) \to (\mathbb{R}, \mathcal{B}(\mathbb{R}))$
ölçülebilir bir fonksiyondur; bu $\{X \leq c\} \in \mathcal{F}$ koşuluyla
eşdeğerdir. Rasgele değişkenler cebirsel işlemler ve limit altında kapalıdır.
\end{mikroozet}

% =====================================================================
\section{Dağılım Fonksiyonu}
\label{sec:dagilim_fonksiyonu}
% =====================================================================

\begin{kesif}
$X$ bir rasgele değişken. $P(X \leq 1.5)$, $P(X \leq 3.7)$ gibi değerleri
biliyorsak $X$ hakkında ne kadar bilgi sahibi oluruz?
Eğer tüm $c \in \mathbb{R}$ için $P(X \leq c)$ değerlerini biliyorsak,
bu yeterli midir — yani bu değerler $X$'in dağılımını tamamen belirler mi?
\end{kesif}

\begin{tanim}[Kümülatif Dağılım Fonksiyonu]
\label{def:kdf}
$X$ bir rasgele değişken olsun. \textbf{Kümülatif dağılım fonksiyonu} (KDF)
(cumulative distribution function, CDF)
\[
  F_X(x) = P(X \leq x), \qquad x \in \mathbb{R}
\]
ile tanımlanır.
\end{tanim}

\begin{teorem}[KDF'nin Karakteristik Özellikleri]
\label{thm:kdf_ozellik}
Her KDF aşağıdaki özellikleri sağlar:
\begin{enumerate}[(i)]
  \item \textbf{Monoton artmayan değil:} $x \leq y \Rightarrow F(x) \leq F(y)$.
  \item \textbf{Sağdan sürekli:} $\lim_{y \downarrow x} F(y) = F(x)$.
  \item \textbf{Sınır davranışı:} $\lim_{x \to -\infty} F(x) = 0$ ve
        $\lim_{x \to +\infty} F(x) = 1$.
\end{enumerate}
Tersine, bu üç özelliği sağlayan her $F : \mathbb{R} \to [0,1]$ fonksiyonu
bir rasgele değişkenin KDF'sidir.
\end{teorem}

\begin{proof}
\textit{(i):} $x \leq y$ ise $\{X \leq x\} \subseteq \{X \leq y\}$, dolayısıyla
$F(x) = P(X \leq x) \leq P(X \leq y) = F(y)$.

\textit{(ii):} $y_n \downarrow x$ olsun. O zaman $\{X \leq y_n\} \downarrow \{X \leq x\}$
(azalan olaylar dizisi, $\bigcap_n \{X \leq y_n\} = \{X \leq x\}$).
Ölçünün yukarıdan sürekliliği (Bölüm~\ref{chp:olcum_teorisi}, Teorem~\ref{thm:olcuozellik}):
\[
  \lim_{y \downarrow x} F(y) = \lim_{n \to \infty} P(X \leq y_n) = P\!\left(\bigcap_n \{X \leq y_n\}\right) = P(X \leq x) = F(x).
\]

\textit{(iii):} $x_n \to -\infty$ ise $\{X \leq x_n\} \downarrow \emptyset$, dolayısıyla
$F(x_n) \to P(\emptyset) = 0$. Benzer şekilde $x_n \to +\infty$ ise
$\{X \leq x_n\} \uparrow \Omega$, dolayısıyla $F(x_n) \to P(\Omega) = 1$.

\textit{Tersine:} $(i)$--$(iii)$'ü sağlayan $F$ verildiğinde,
$\mathbb{R}$ üzerindeki Lebesgue-Stieltjes ölçüsü inşası yoluyla
(Carathéodory genişleme teoremi) $P((a,b]) = F(b) - F(a)$ sağlayan
bir olasılık ölçüsü $P$ inşa edilebilir; bu ölçünün üreteceği KDF tam olarak $F$'dir.
\end{proof}

\begin{onerme}[KDF'den Olasılık Hesabı]
\label{prop:kdf_hesap}
$F = F_X$ KDF'si olan $X$ için:
\begin{align*}
  P(X > x)   &= 1 - F(x), \\
  P(a < X \leq b) &= F(b) - F(a), \\
  P(X = x)   &= F(x) - F(x^-) \quad \text{(sol limit: } F(x^-) = \lim_{y \uparrow x} F(y)\text{)}, \\
  P(a \leq X \leq b) &= F(b) - F(a^-), \\
  P(a < X < b) &= F(b^-) - F(a).
\end{align*}
\end{onerme}

\begin{proof}
$P(a < X \leq b) = P(X \leq b) - P(X \leq a) = F(b) - F(a)$.
$P(X = x) = P(X \leq x) - P(X < x) = F(x) - \lim_{y \uparrow x} F(y) = F(x) - F(x^-)$.
Diğerleri bu iki formülden türetilir.
\end{proof}

\begin{hataanalizi}
"$P(X < x) = F(x)$" hatası: Sol açık aralık için $P(X < x) = F(x^-)$
olup $F$'nin $x$'te süreksiz olduğu noktalarda $F(x^-) < F(x)$.
Sağdan sürekli tanım $P(X \leq x) = F(x)$'i verir; $P(X < x)$ farklıdır.
\end{hataanalizi}

\begin{kendinleifade}
KDF'nin üç temel özelliğini (monoton, sağdan sürekli, sınırlar)
bir yabancı dil örneğiyle canlandırın: örneğin "dil testi skoru"nun
KDF'si nasıl görünürdü? Sürekli mi olurdu, yoksa sıçramalı mı?
\ogrencialani[2.5cm]
\end{kendinleifade}

\begin{mikroozet}
$F_X(x) = P(X \leq x)$. KDF: monoton, sağdan sürekli, $0 \to 1$ sınırları.
$P(X = x) = F(x) - F(x^-)$; süreklilik noktalarında bu sıfırdır.
\end{mikroozet}

% =====================================================================
\section{Kesikli Dağılımlar}
\label{sec:kesikli}
% =====================================================================

\begin{tanim}[Kesikli Rasgele Değişken]
\label{def:kesikli}
$X$ rasgele değişkeni \textbf{kesikli} (discrete) dır, eğer
sayılabilir bir $S \subset \mathbb{R}$ kümesi için $P(X \in S) = 1$ ise.
$S$'nin elemanlarına $X$'in \textbf{olası değerleri} denir.
\textbf{Olasılık kütle fonksiyonu} (probability mass function, PMF):
\[
  p_X(x) = P(X = x), \qquad x \in S.
\]
\end{tanim}

\begin{teorem}[PMF-KDF İlişkisi]
\label{thm:pmf_kdf}
$X$ kesikli rasgele değişken, $S = \{x_1, x_2, \ldots\}$ ($x_1 < x_2 < \cdots$) olsun. O zaman
\[
  F_X(x) = \sum_{x_i \leq x} p_X(x_i)
\]
ve $F_X$, her $x_i$'de büyüklüğü $p_X(x_i)$ olan sıçramayla
sağdan sürekli bir basamak fonksiyonudur.
\end{teorem}

\begin{proof}
$\{X \leq x\} = \bigsqcup_{x_i \leq x} \{X = x_i\}$ ayrık birleşimidir.
Sayılabilir toplamlılık:
$F_X(x) = P(X \leq x) = \sum_{x_i \leq x} P(X = x_i) = \sum_{x_i \leq x} p_X(x_i)$.
$y \downarrow x$ alındığında $\{x_i \leq y\}$ kümesi $y$'nin küçük aralıklarda
değişimi boyunca $\{x_i \leq x\}$'e yakınsadığından sağdan süreklilik sağlanır.
\end{proof}

\begin{ornek}[Binom]
\label{ex:binom_pmf}
$X \sim \Binom(n, p)$: $n$ bağımsız Bernoulli denemesinde başarı sayısı.
\[
  p_X(k) = \binom{n}{k} p^k (1-p)^{n-k}, \quad k = 0, 1, \ldots, n.
\]
$\sum_{k=0}^n p_X(k) = (p + (1-p))^n = 1$ (Binom teoremi).
\end{ornek}

% =====================================================================
\section{Mutlak Sürekli Dağılımlar ve Yoğunluk Fonksiyonu}
\label{sec:surekli}
% =====================================================================

\begin{kesif}
Bir terazi ibresi [0,1] aralığında rasgele bir değere işaret ediyor.
$P(X = 0.5) = ?$ — sonsuz çok nokta var, her birinin olasılığı sıfır olmalı.
Peki $P(0.3 \leq X \leq 0.7) = ?$ Bu sıfır değil.
Nasıl bir kavramla olasılıkları hesaplayabiliriz?
\end{kesif}

\begin{tanimgerekce}
Radon-Nikodym teoremi der ki: $P_X \ll \lambda$ (Lebesgue ölçüsüne göre
mutlak sürekli) ise $P_X(B) = \int_B f(x) \, d\lambda(x)$ yazan bir $f \geq 0$ var.
Bu $f$ "yoğunluk"tur — noktasal olasılık değil, uzunluk başına olasılık yoğunluğu.
\end{tanimgerekce}

\begin{tanim}[Mutlak Sürekli Rasgele Değişken]
\label{def:surekli_rd}
$X$ rasgele değişkenin \textbf{olasılık yoğunluk fonksiyonu} (probability density function, PDF)
$f_X : \mathbb{R} \to [0, \infty)$ mevcutsa ve her $B \in \mathcal{B}(\mathbb{R})$ için
\[
  P(X \in B) = \int_B f_X(x) \, dx
\]
sağlanıyorsa $X$'e \textbf{mutlak sürekli} (absolutely continuous) rasgele değişken denir.
Özellikle $F_X(x) = \int_{-\infty}^x f_X(t) \, dt$ ve $\int_{-\infty}^{+\infty} f_X(x) \, dx = 1$.
\end{tanim}

\begin{teorem}[PDF-KDF İlişkisi]
\label{thm:pdf_kdf}
$X$ mutlak sürekli rasgele değişken ve $f_X$ PDF'si olsun. O zaman:
\begin{enumerate}[(i)]
  \item $F_X(x) = \int_{-\infty}^x f_X(t) \, dt$, dolayısıyla $F_X$ süreklidir.
  \item $f_X$ sürekli olduğu her $x$ noktasında $F_X'(x) = f_X(x)$.
  \item Her $x$ için $P(X = x) = 0$.
  \item $f_X \geq 0$ n.k.\ ve $\int_{\mathbb{R}} f_X = 1$.
\end{enumerate}
\end{teorem}

\begin{proof}
\textit{(i)}: Tanımdan; $(-\infty, x] \in \mathcal{B}(\mathbb{R})$ seçilir.
\textit{(ii)}: Lebesgue farklılaştırma teoremi; $f_X$ sürekli ise
$\frac{d}{dx} \int_{-\infty}^x f_X(t) \, dt = f_X(x)$.
\textit{(iii)}: $P(X = x) = F_X(x) - F_X(x^-) = 0$ ($F_X$ sürekli olduğundan).
\textit{(iv)}: $f_X \geq 0$ a.e., $\int f_X = P(X \in \mathbb{R}) = 1$.
\end{proof}

\begin{hataanalizi}
"$f_X(x) = P(X = x)$" en yaygın hatadır. Doğrusu: $f_X(x) \, dx \approx P(x < X \leq x + dx)$
yani $f_X(x)$, küçük aralığa düşen \emph{olasılık yoğunluğudur}, bizzat olasılık değildir.
$f_X(x) > 1$ olabilir (örneğin $X \sim \mathrm{Uniform}(0, 0.1)$ ise $f_X(x) = 10$).
\end{hataanalizi}

\begin{karsiornek}
$f_X(x) = 10 \cdot \mathbf{1}_{[0, 0.1]}(x)$: $f_X(x) = 10 > 1$ ama yine de geçerli bir PDF.
$\int_0^{0.1} 10 \, dx = 1$. "Yoğunluk 1'i geçemez" iddiası yanlıştır.
\end{karsiornek}

\begin{ornek}[Normal Dağılım]
$X \sim \Normal(\mu, \sigma^2)$: PDF
\[
  f_X(x) = \frac{1}{\sigma\sqrt{2\pi}} \exp\!\left(-\frac{(x-\mu)^2}{2\sigma^2}\right).
\]
$\int_{-\infty}^{+\infty} f_X(x) \, dx = 1$ doğrulaması: Bölüm~\ref{chp:olcum_teorisi}'deki
Gauss integrali hesabına bakınız (Fubini-Tonelli).
\end{ornek}

\begin{mikroozet}
PDF: $P(X \in B) = \int_B f_X \, dx$; $f_X \geq 0$, $\int f_X = 1$.
$f_X(x)$ bir olasılık değil, yoğunluktur; $f_X > 1$ olabilir.
Her tekil noktanın olasılığı sıfırdır.
\end{mikroozet}

% =====================================================================
\section{Karma Dağılımlar ve Lebesgue Ayrışımı}
\label{sec:karma}
% =====================================================================

\begin{tanim}[Karma Dağılım]
\label{def:karma}
$X$ ne kesikli ne de mutlak sürekli olan bir rasgele değişken ise,
KDF'si $F_X = p \cdot F_{\text{kes}} + (1-p) \cdot F_{\text{sür}}$
biçiminde ($0 < p < 1$) \textbf{karma dağılım} (mixed distribution) olarak
yazılabilir; burada $F_{\text{kes}}$ kesikli, $F_{\text{sür}}$ mutlak sürekli
bir KDF'dir.
\end{tanim}

\begin{onerme}[Lebesgue Ayrışım Teoremi — Özel Durum]
\label{prop:lebesgue_ayrisim}
Her $F : \mathbb{R} \to [0,1]$ KDF'si şu biçimde tek türlü yazılabilir:
\[
  F = \alpha_1 F_{\text{saf-kes}} + \alpha_2 F_{\text{mut-sür}} + \alpha_3 F_{\text{tekil-sür}},
\]
$\alpha_1 + \alpha_2 + \alpha_3 = 1$, $\alpha_i \geq 0$; burada
$F_{\text{tekil-sür}}$ sürekli ama türevi neredeyse her yerde sıfır olan
(Lebesgue ölçüsüne göre mutlak sürekli olmayan) bir KDF'dir.
Cantor fonksiyonu (şeytan merdiveni) $F_{\text{tekil-sür}}$'ye örnektir.
\end{onerme}

\begin{hocanotu}
Pratikte tekil-sürekli bileşen nadiren karşımıza çıkar; uygulamalar
büyük ölçüde $\alpha_3 = 0$ durumunda (karma = kesikli + mutlak sürekli) çalışır.
Ancak tam teorik çerçeve için bu bileşeni bilmek önemlidir.
\end{hocanotu}

\begin{ornek}[Karma Dağılım]
Bir sigorta poliçesinde:
$P(X = 0) = 0.3$ (hasar yok), $X > 0$ durumunda $X \sim \mathrm{Exp}(\lambda)$.
KDF: $F_X(x) = 0$ ($x < 0$); $F_X(0) = 0.3$; $F_X(x) = 0.3 + 0.7(1 - e^{-\lambda x})$ ($x > 0$).
Ne tamamen kesikli (sürekli bir parçası var) ne tamamen sürekli (0'da sıçrama var).
\end{ornek}

% =====================================================================
\section{Birden Fazla Rasgele Değişken: Ortak Dağılımlar}
\label{sec:rd_ortak}
% =====================================================================

\begin{tanim}[Rasgele Vektör ve Ortak KDF]
\label{def:ortak_kdf}
$X_1, \ldots, X_n : \Omega \to \mathbb{R}$ rasgele değişkenler olsun.
$(X_1, \ldots, X_n) : \Omega \to \mathbb{R}^n$ \textbf{rasgele vektördür}
(ölçülebilirlik: $\mathcal{F}$'den $\mathcal{B}(\mathbb{R}^n)$'ye).
\textbf{Ortak KDF} (joint CDF):
\[
  F_{X_1,\ldots,X_n}(x_1, \ldots, x_n) = P(X_1 \leq x_1, \ldots, X_n \leq x_n).
\]
İki değişkenli durumda:
\begin{itemize}
  \item \textbf{Ortak PMF} (kesikli): $p(x,y) = P(X=x, Y=y)$.
  \item \textbf{Ortak PDF} (mutlak sürekli): $f_{X,Y}(x,y) \geq 0$,
        $P((X,Y) \in B) = \iint_B f_{X,Y}(x,y) \, dx \, dy$.
\end{itemize}
\end{tanim}

% =====================================================================
\section{Marjinal ve Koşullu Dağılımlar}
\label{sec:marjinal_kosullu}
% =====================================================================

\begin{tanim}[Marjinal Dağılım]
\label{def:marjinal}
$(X, Y)$ ortak dağılımlı rasgele vektör olsun.
$X$'in \textbf{marjinal dağılımı}:
\begin{itemize}
  \item Kesikli: $p_X(x) = \sum_y p(x, y)$.
  \item Sürekli: $f_X(x) = \int_{-\infty}^{+\infty} f_{X,Y}(x,y) \, dy$.
\end{itemize}
\end{tanim}

\begin{kritiksorgu}
"Aynı marjinallere sahip iki rasgele vektörün ortak dağılımı aynı mıdır?"
Hayır: Marjinaller ortak dağılımı belirlemez.
Örnek: $X \sim \mathrm{Uniform}(-1,1)$, $Y = X$ ile $Y = -X$ seçimleri
aynı marjinalleri verir ama farklı ortak dağılımlar üretir.
Bu gerçek, istatistiksel bağımlılık teorisinin (kopula, korelasyon) özüdür.
\end{kritiksorgu}

\begin{tanim}[Koşullu Dağılım]
\label{def:kosullu_dagilim}
\textbf{Kesikli durum:} $p_Y(y) > 0$ ise
\[
  p_{X|Y}(x \mid y) = P(X = x \mid Y = y) = \frac{p(x,y)}{p_Y(y)}.
\]
\textbf{Sürekli durum:} $f_Y(y) > 0$ ise
\[
  f_{X|Y}(x \mid y) = \frac{f_{X,Y}(x,y)}{f_Y(y)}.
\]
$f_{X|Y}(\cdot \mid y)$, $y$ sabit tutulduğunda $x$'in bir PDF'sidir:
$\int f_{X|Y}(x \mid y) \, dx = 1$.
\end{tanim}

\begin{teorem}[Toplam Olasılık — Sürekli Versiyon]
\label{thm:toplam_surekli}
$(X, Y)$ ortak sürekli dağılımlı olsun. O zaman
\[
  f_X(x) = \int_{-\infty}^{+\infty} f_{X|Y}(x \mid y) \, f_Y(y) \, dy.
\]
\end{teorem}

\begin{proof}
$f_X(x) = \int f_{X,Y}(x,y) \, dy = \int \frac{f_{X,Y}(x,y)}{f_Y(y)} \cdot f_Y(y) \, dy
= \int f_{X|Y}(x \mid y) \, f_Y(y) \, dy$.
\end{proof}

\begin{ornek}[Koşullu Dağılım Hesabı]
$f_{X,Y}(x,y) = 6x$ ($0 < x < y < 1$), sıfır başka yerde.
Marjinal: $f_Y(y) = \int_0^y 6x \, dx = 3y^2$.
Koşullu: $f_{X|Y}(x \mid y) = \frac{6x}{3y^2} = \frac{2x}{y^2}$, $0 < x < y$.
\end{ornek}

% =====================================================================
\section{Bağımsız Rasgele Değişkenler}
\label{sec:rd_bagimsizlik}
% =====================================================================

\begin{tanim}[Rasgele Değişken Bağımsızlığı]
\label{def:rd_bagimsiz}
$X_1, \ldots, X_n$ rasgele değişkenler \textbf{bağımsızdır} (independent)
ancak ve ancak her $B_1, \ldots, B_n \in \mathcal{B}(\mathbb{R})$ için
\[
  P(X_1 \in B_1, \ldots, X_n \in B_n) = \prod_{i=1}^n P(X_i \in B_i).
\]
Eşdeğer olarak $\sigma(X_1), \ldots, \sigma(X_n)$ $\sigma$-cebirleri bağımsızdır.
\end{tanim}

\begin{teorem}[Bağımsızlığın Eşdeğer Koşulları]
\label{thm:bagimsiz_denk}
$(X, Y)$ rasgele vektör için aşağıdakiler eşdeğerdir:
\begin{enumerate}[(i)]
  \item $X$ ve $Y$ bağımsızdır.
  \item Her $x, y$ için $F_{X,Y}(x,y) = F_X(x) \cdot F_Y(y)$.
  \item \textit{(Kesikli:)} Her $x, y$ için $p(x,y) = p_X(x) \cdot p_Y(y)$.
  \item \textit{(Sürekli:)} $f_{X,Y}(x,y) = f_X(x) \cdot f_Y(y)$ n.k.
  \item $f_{X|Y}(x \mid y) = f_X(x)$ n.k.\ (koşullu dağılım marjinale eşittir).
\end{enumerate}
\end{teorem}

\begin{proof}
$(i) \Leftrightarrow (ii)$: $B_i = (-\infty, x_i]$ seçilirse $(i)$'den $(ii)$ çıkar;
$\pi$-sistemi argümanıyla tersi de kurulur.
$(ii) \Leftrightarrow (iv)$: Sürekli durumda $F_{X,Y}(x,y) = \int_{-\infty}^x \int_{-\infty}^y f(s,t) \, dt \, ds$;
bu $F_X(x) F_Y(y)$'ye eşit olması $f = f_X \cdot f_Y$ a.e.\ ile eşdeğerdir
(Fubini).
$(iv) \Leftrightarrow (v)$: $f_{X|Y}(x \mid y) = f(x,y)/f_Y(y) = f_X(x)$.
\end{proof}

\begin{hataanalizi}
"Korelasyonsuz $\Rightarrow$ bağımsız" hatası: $\mathrm{Cor}(X, Y) = 0$
bağımsızlık için yeterli değildir.
Örnek: $X \sim \mathrm{Uniform}(-1,1)$, $Y = X^2$. O zaman $E[XY] = E[X^3] = 0 = E[X]E[Y]$,
yani $\mathrm{Cov}(X, Y) = 0$, ama $Y$ tamamen $X$'e bağımlıdır.
Yalnızca normal dağılım için $\mathrm{Cor} = 0 \Leftrightarrow$ bağımsızlık.
\end{hataanalizi}

\begin{mikroozet}
Bağımsızlık: $F_{X,Y} = F_X \cdot F_Y$ (veya PDF çarpımlanması).
Bağımsızlık $\Rightarrow$ korelasyonsuzluk; tersi genel olarak yanlış.
\end{mikroozet}

% =====================================================================
\section{Rasgele Değişkenlerin Fonksiyonları}
\label{sec:rd_donusum}
% =====================================================================

\begin{motivasyon}[Neden Dönüşüm?]
$X$ bilinen dağılımlı ise $Y = g(X)$'in dağılımı nedir?
Bu soru istatistiksel hesaplarda sürekli karşımıza çıkar:
örneğin $X \sim \Normal(0,1)$ ise $Y = X^2$'nin dağılımı chi-kare,
$Y = e^X$ ise lognormal dağılım verir.
\end{motivasyon}

\subsection{Tekil Boyutlu Dönüşüm}
\label{subsec:donusum_1d}

\begin{teorem}[Birebir Dönüşüm Yöntemi]
\label{thm:donusum_1d}
$X$ sürekli rasgele değişken, PDF'si $f_X$. $g : \mathbb{R} \to \mathbb{R}$
monoton, türevlenebilir ve birebir olsun, $Y = g(X)$.
$h = g^{-1}$ (ters fonksiyon) olmak üzere:
\[
  f_Y(y) = f_X(h(y)) \cdot \left|\frac{dh}{dy}\right|.
\]
\end{teorem}

\begin{proof}
\textbf{$g$ artan} olsun. O zaman
\[
  F_Y(y) = P(Y \leq y) = P(g(X) \leq y) = P(X \leq g^{-1}(y)) = F_X(h(y)).
\]
Türev alınca: $f_Y(y) = f_X(h(y)) \cdot h'(y) = f_X(h(y)) \cdot |h'(y)|$.

\textbf{$g$ azalan} olsun. O zaman
\[
  F_Y(y) = P(Y \leq y) = P(g(X) \leq y) = P(X \geq h(y)) = 1 - F_X(h(y)).
\]
Türev: $f_Y(y) = -f_X(h(y)) \cdot h'(y) = f_X(h(y)) \cdot |h'(y)|$ ($h' < 0$ olduğundan).
\end{proof}

\begin{ornek}[Lognormal]
$X \sim \Normal(\mu, \sigma^2)$, $Y = e^X$ (lognormal dağılım).
$h(y) = \ln y$, $|h'(y)| = 1/y$.
\[
  f_Y(y) = \frac{1}{y\sigma\sqrt{2\pi}} \exp\!\left(-\frac{(\ln y - \mu)^2}{2\sigma^2}\right), \quad y > 0.
\]
\end{ornek}

\begin{onerme}[Genel Monoton Olmayan Dönüşüm]
\label{prop:genel_donusum}
$g$ birebir olmayabilir. $g^{-1}(y) = \{x : g(x) = y\}$ sonlu küme ise:
\[
  f_Y(y) = \sum_{x : g(x) = y} f_X(x) \cdot \left|\frac{dg}{dx}\bigg|_{x}\right|^{-1}.
\]
\end{onerme}

\begin{ornek}[$Y = X^2$, $X \sim \Normal(0,1)$]
$g(x) = x^2$, $g^{-1}(y) = \{\sqrt{y}, -\sqrt{y}\}$ ($y > 0$).
$|dg/dx|^{-1} = 1/(2|x|)$.
\[
  f_Y(y) = f_X(\sqrt{y}) \cdot \frac{1}{2\sqrt{y}} + f_X(-\sqrt{y}) \cdot \frac{1}{2\sqrt{y}}
         = \frac{1}{\sqrt{2\pi y}} e^{-y/2}, \quad y > 0.
\]
Bu chi-kare dağılımının ($\chi^2_1$) PDF'sidir.
\end{ornek}

\subsection{Çok Boyutlu Dönüşüm ve Jacobian}
\label{subsec:donusum_jacobian}

\begin{teorem}[Jacobian Formülü]
\label{thm:jacobian}
$(X_1, \ldots, X_n)$ ortak sürekli dağılımlı; PDF'si $f_{X_1,\ldots,X_n}$.
$g : \mathbb{R}^n \to \mathbb{R}^n$ birebir ve türevlenebilir olsun,
$(Y_1, \ldots, Y_n) = g(X_1, \ldots, X_n)$.
Ters dönüşüm $h = g^{-1}$ ve Jacobian matrisi
\[
  J = \frac{\partial(x_1,\ldots,x_n)}{\partial(y_1,\ldots,y_n)}
    = \left(\frac{\partial h_i}{\partial y_j}\right)_{n \times n}
\]
olmak üzere:
\[
  f_{Y_1,\ldots,Y_n}(y_1,\ldots,y_n) = f_{X_1,\ldots,X_n}(h_1(y),\ldots,h_n(y)) \cdot |\det J|.
\]
\end{teorem}

\begin{proof}
Çok değişkenli değişken değiştirme formülü:
$P((Y_1,\ldots,Y_n) \in B) = \int_B f_Y \, dy = \int_{g^{-1}(B)} f_X \, dx$.
$dx = |\det J| \, dy$ (Jacobian dönüşümü); ikisini eşitleyince $f_Y$ elde edilir.
\end{proof}

\begin{ornek}[Polar Dönüşüm]
$(X, Y)$ bağımsız, her ikisi $\Normal(0,1)$.
$R = \sqrt{X^2+Y^2}$, $\Theta = \arctan(Y/X)$.
Ters: $x = r\cos\theta$, $y = r\sin\theta$.
$\det J = r$. Ortak PDF $f_{R,\Theta}(r, \theta) = f_{X,Y}(r\cos\theta, r\sin\theta) \cdot r$.
\[
  f_{X,Y}(x,y) = \frac{1}{2\pi} e^{-(x^2+y^2)/2}
  \implies f_{R,\Theta}(r,\theta) = r e^{-r^2/2} \cdot \frac{1}{2\pi},
\]
$r > 0$, $\theta \in [0, 2\pi)$.
Marjinaller: $\Theta \sim \mathrm{Uniform}(0, 2\pi)$ ve $R^2 \sim \chi^2_2$ (Rayleigh dağılımı).
Bu Box-Muller dönüşümünün temelidir.
\end{ornek}

\subsection{Toplam Dağılımı: Evrişim}
\label{subsec:evrisim}

\begin{tanim}[Evrişim (Konvolüsyon)]
\label{def:evrisim}
$X, Y$ bağımsız, ortak sürekli dağılımlı; $Z = X + Y$ olsun. O zaman
\[
  f_Z(z) = (f_X * f_Y)(z) = \int_{-\infty}^{+\infty} f_X(z - y) \, f_Y(y) \, dy.
\]
Kesikli durumda: $p_Z(z) = \sum_y p_X(z-y) \, p_Y(y)$.
\end{tanim}

\begin{proof}
$F_Z(z) = P(X + Y \leq z) = \int_{-\infty}^{+\infty} P(X \leq z - y) \, f_Y(y) \, dy
= \int_{-\infty}^{+\infty} F_X(z-y) \, f_Y(y) \, dy$.
$z$'ye göre türev: Leibniz kuralıyla $f_Z(z) = \int f_X(z-y) f_Y(y) \, dy$.
\end{proof}

\begin{ornek}[Normal + Normal = Normal]
$X \sim \Normal(\mu_1, \sigma_1^2)$, $Y \sim \Normal(\mu_2, \sigma_2^2)$ bağımsız.
Evrişim hesabı: $Z = X + Y \sim \Normal(\mu_1+\mu_2, \sigma_1^2+\sigma_2^2)$.
Hesap, Gauss integrallerini tamamlama yöntemiyle yapılır;
Bölüm~\ref{chp:beklenti}'deki karakteristik fonksiyon yöntemi çok daha hızlıdır.
\end{ornek}

% ---- Disiplinlerarası Bağlantı (TYMM) ----
\begin{kopru}[Disiplinlerarası — Finans ve Mühendislik]
\textbf{Finans:} Hisse senedi getirisinin lognormal modellenmesi (Black-Scholes),
rasgele değişken dönüşüm teorisinin doğrudan uygulamasıdır.\\
\textbf{Mühendislik:} Haberleşmede sinyal+gürültü modeli $Z = S + N$,
$N \sim \Normal(0, \sigma^2)$: evrişim yoluyla alıcı tasarımı.\\
\textbf{Fizik:} Kuantum mekaniğinde konum ve momentum dağılımları
PDF kavramıyla doğrudan ifade edilir; Fourier dönüşümü ile karakteristik
fonksiyon arasındaki bağlantı Heisenberg belirsizlik ilkesini yansıtır.
\end{kopru}

% ---- Öz-Yansıtma ----
\begin{ozyansima}
\begin{itemize}[leftmargin=1.8em]
  \item "Rasgele değişken bir sayı mı, yoksa fonksiyon mu?" — cevabınız ve gerekçeniz?
  \item KDF'nin hangi özelliği sizi en çok şaşırttı?
  \item Sürekli durumda $P(X = x) = 0$ gerçeği sezgiye aykırı geliyor mu?
        Bunu fizik veya mühendislikten bir örnekle nasıl uzlaştırırsınız?
  \item Bağımsızlık ile korelasyonsuzluk arasındaki fark neden önemlidir?
\end{itemize}
\end{ozyansima}

% ---- Köprü Notu ----
\begin{kopru}
\textbf{Önceki bölüm:} Bölüm~\ref{chp:olasilik_uzaylari} — Kolmogorov aksiyomları,
bağımsızlık olaylar düzeyinde; burada rasgele değişkenlere taşındı.\\
\textbf{Sonraki bölüm:} Bölüm~\ref{chp:beklenti} — Dağılımlar verildiğinde
$E[X]$, $\mathrm{Var}(X)$, momentler; Lebesgue integrali olarak beklenti.\\
\textbf{Dijital katman:} Dönüşüm formüllerinin sayısal doğrulaması ve
Monte Carlo simülasyonu \texttt{bolum03\_donusumler.ipynb} defterinde.
\defterbaglantisi{bolum03\_donusumler.ipynb}{%
  Jacobian dönüşümü, evrişim ve KDF simülasyonu}
\end{kopru}

% =====================================================================
\section{Çok Boyutlu Normal Dağılım}
\label{sec:cok_boyutlu_normal}
% =====================================================================

\begin{motivasyon}[Neden Çok Boyutlu Normal?]
Tek boyutlu Normal dağılım, pek çok bağımsız etkinin toplamının limit dağılımı
(MLT) olarak doğal biçimde ortaya çıkar.
Çok boyutlu uzantısı, bağımlı Normal rasgele değişkenlerin tam ailesini verir
ve lineer dönüşümler altında kapalıdır.
İstatistik, makine öğrenmesi ve sinyal işlemede varsayılan "default" dağılımdır.
\end{motivasyon}

\begin{tanim}[Çok Boyutlu Normal Dağılım]
\label{def:cok_boyutlu_normal}
$\mathbf{X} = (X_1, \ldots, X_k)^\top$ rasgele vektör.
$\mathbf{X} \sim \Normal_k(\boldsymbol{\mu}, \boldsymbol{\Sigma})$
(\textbf{$k$-boyutlu Normal}, multivariate normal) denir eğer:
\begin{enumerate}[(i)]
  \item $\boldsymbol{\mu} = E[\mathbf{X}] \in \mathbb{R}^k$ ortalama vektör.
  \item $\boldsymbol{\Sigma} = \mathrm{Cov}(\mathbf{X}) \in \mathbb{R}^{k\times k}$ simetrik pozitif yarı-belirli kovaryans matrisi.
  \item $\boldsymbol{\Sigma}$ pozitif belirli ise ortak PDF:
  \[
    f(\mathbf{x}) = \frac{1}{(2\pi)^{k/2}|\boldsymbol{\Sigma}|^{1/2}}
    \exp\!\left(-\frac{1}{2}(\mathbf{x}-\boldsymbol{\mu})^\top \boldsymbol{\Sigma}^{-1}(\mathbf{x}-\boldsymbol{\mu})\right).
  \]
\end{enumerate}
Eşdeğer tanım: $\mathbf{X} = \boldsymbol{\mu} + A\mathbf{Z}$, $\mathbf{Z} \sim \Normal_m(\mathbf{0}, I_m)$ bağımsız,
$\boldsymbol{\Sigma} = AA^\top$.
\end{tanim}

\begin{teorem}[Çok Boyutlu Normal'in Temel Özellikleri]
\label{thm:cb_normal}
$\mathbf{X} \sim \Normal_k(\boldsymbol{\mu}, \boldsymbol{\Sigma})$:
\begin{enumerate}[(i)]
  \item \textbf{Marjinaller:} $X_i \sim \Normal(\mu_i, \sigma_{ii})$ (marjinallar normal).
  \item \textbf{Lineer dönüşüm:} $A\mathbf{X} + \mathbf{b} \sim \Normal(A\boldsymbol{\mu}+\mathbf{b}, A\boldsymbol{\Sigma}A^\top)$.
  \item \textbf{Bağımsızlık:} $X_i$ ve $X_j$ bağımsız $\Leftrightarrow$ $\sigma_{ij} = 0$ (Normal için korelasyonsuzluk = bağımsızlık!).
  \item \textbf{Koşullu dağılım:} $\mathbf{X}$'i $(\mathbf{X}_1, \mathbf{X}_2)$ şeklinde parçalayın.
  $\mathbf{X}_1 \mid \mathbf{X}_2 = \mathbf{x}_2 \sim \Normal(\boldsymbol{\mu}_{1|2},\, \boldsymbol{\Sigma}_{11|2})$:
  \begin{align*}
    \boldsymbol{\mu}_{1|2} &= \boldsymbol{\mu}_1 + \boldsymbol{\Sigma}_{12}\boldsymbol{\Sigma}_{22}^{-1}(\mathbf{x}_2 - \boldsymbol{\mu}_2),\\
    \boldsymbol{\Sigma}_{11|2} &= \boldsymbol{\Sigma}_{11} - \boldsymbol{\Sigma}_{12}\boldsymbol{\Sigma}_{22}^{-1}\boldsymbol{\Sigma}_{21}.
  \end{align*}
\end{enumerate}
\end{teorem}

\begin{proof}[(i) ve (ii)]
(i) Marjinal: $X_i = \mathbf{e}_i^\top \mathbf{X}$ ($\mathbf{e}_i$ birim vektör); (ii) özel durumu.\\
(ii) KF: $\varphi_{A\mathbf{X}+\mathbf{b}}(\mathbf{t}) = e^{i\mathbf{t}^\top \mathbf{b}} \varphi_{\mathbf{X}}(A^\top \mathbf{t})$.
$\varphi_{\mathbf{X}}(\mathbf{t}) = e^{i\mathbf{t}^\top\boldsymbol{\mu} - \mathbf{t}^\top\boldsymbol{\Sigma}\mathbf{t}/2}$ (Normal KF).
$\varphi_{A\mathbf{X}+\mathbf{b}}(\mathbf{t}) = e^{i\mathbf{t}^\top(A\boldsymbol{\mu}+\mathbf{b}) - \mathbf{t}^\top A\boldsymbol{\Sigma}A^\top\mathbf{t}/2}$: $\Normal(A\boldsymbol{\mu}+\mathbf{b}, A\boldsymbol{\Sigma}A^\top)$ KF'si.
\end{proof}

\begin{ornek}[İki Boyutlu Normal]
$k=2$: $\boldsymbol{\mu} = (\mu_1, \mu_2)^\top$,
$\boldsymbol{\Sigma} = \begin{pmatrix}\sigma_1^2 & \rho\sigma_1\sigma_2 \\ \rho\sigma_1\sigma_2 & \sigma_2^2\end{pmatrix}$.
PDF:
\[
  f(x_1,x_2) = \frac{1}{2\pi\sigma_1\sigma_2\sqrt{1-\rho^2}}
  \exp\!\left(-\frac{1}{2(1-\rho^2)}\left[\frac{(x_1-\mu_1)^2}{\sigma_1^2}
  - \frac{2\rho(x_1-\mu_1)(x_2-\mu_2)}{\sigma_1\sigma_2}
  + \frac{(x_2-\mu_2)^2}{\sigma_2^2}\right]\right).
\]
$X_1 \mid X_2 = x_2 \sim \Normal\!\left(\mu_1 + \rho\frac{\sigma_1}{\sigma_2}(x_2-\mu_2),\,
\sigma_1^2(1-\rho^2)\right)$.
\end{ornek}

\begin{hataanalizi}
"Marjinalleri Normal olan rasgele vektör çok boyutlu Normal'dir."
Yanlış: Korrelasyon yapısı yanlış olabilir. Örnek: $Z \sim \Normal(0,1)$,
$W = Z \cdot \mathrm{sgn}(Z)$; hem $Z$ hem $W$ normal marjinale sahip
ama $(Z, W)$ çift boyutlu normal değil (ortak PDF oval konturlar çizmiyor).
Çok boyutlu Normal tanımı KF veya affine dönüşüm koşuluyla yapılmalı.
\end{hataanalizi}

\begin{mikroozet}
$\mathbf{X}\sim\Normal_k(\boldsymbol{\mu},\boldsymbol{\Sigma})$: marjinaller normal; lineer dönüşümler kapalı;
Normal'de korelasyonsuzluk = bağımsızlık; koşullu dağılım da normaldir.
\end{mikroozet}

% =====================================================================
\section{Kuantil Fonksiyon ve Olasılık İntegral Dönüşümü}
\label{sec:kuantil_fonksiyon}
% =====================================================================

\begin{tanim}[Kuantil Fonksiyon]
\label{def:kuantil_fonksiyon}
$F$ KDF için \textbf{kuantil fonksiyon} (quantile function), genelleştirilmiş ters:
\[
  F^{-1}(u) = \inf\{x : F(x) \geq u\}, \quad u \in (0,1).
\]
$u$'nun $p$. yüzdelik (\textit{percentile}): $F^{-1}(p)$.
$F^{-1}(1/2)$: medyan. $F^{-1}(1/4), F^{-1}(3/4)$: çeyrekler.
\end{tanim}

\begin{teorem}[Olasılık İntegral Dönüşümü]
\label{thm:oit}
\textbf{İleri:} $X$ sürekli rasgele değişken, $F_X$ KDF. $F_X(X) \sim \Uniform(0,1)$.\\
\textbf{Geri:} $U \sim \Uniform(0,1)$. $X = F^{-1}(U)$ rasgele değişkeni $F$ KDF'ye sahip.
\end{teorem}

\begin{proof}
\textbf{İleri:} $P(F_X(X) \leq u) = P(X \leq F_X^{-1}(u)) = F_X(F_X^{-1}(u)) = u$.
(Son eşitlik: $F_X$ sürekli ise $F_X \circ F_X^{-1} = \mathrm{id}$.)\\
\textbf{Geri:} $P(F^{-1}(U) \leq x) = P(U \leq F(x)) = F(x)$ ($U$ uniform).
\end{proof}

\begin{ornek}[Simülasyon Uygulaması]
$\Exp(\lambda)$ dağılımından örneklem üretmek:
$F^{-1}(u) = -\ln(1-u)/\lambda$.
$U \sim \Uniform(0,1)$: $X = -\ln(1-U)/\lambda \sim \Exp(\lambda)$.
Bu standart simülasyon tekniğinin (Inverse Transform Method) matematiksel temelidir.
Normal: $\Phi^{-1}(U)$ kapalı form yok; sayısal yöntemler ($\Phi^{-1} \approx$ probit fonksiyon) kullanılır.
\end{ornek}

\begin{teorem}[Kuantil-Kuantil İlişkisi]
\label{thm:qq_iliskisi}
$X \sim F_X$, $Y \sim F_Y$. $Y \overset{d}{=} F_Y^{-1}(F_X(X))$.
\emph{QQ-plot} prensibi: $X$ örneklemi ile $F_Y^{-1}$ karşılaştırılarak $F_Y$ uyumu test edilir.
\end{teorem}

\begin{mikroozet}
$F^{-1}(u) = \inf\{x: F(x)\geq u\}$: kuantil fonksiyon.
OID: $F(X)\sim\Uniform$; $F^{-1}(U)\sim F$.
Simülasyon: her sürekli dağılım $\Uniform(0,1)$'den üretilir.
\end{mikroozet}

% =====================================================================
%  ALIŞTIRMALAR
% =====================================================================

\cozumlualistirmalar

\begin{alistirma}[Al.3.1]{A}{Ölçülebilirlik Kontrolü}
$\Omega = \{1, 2, 3, 4\}$, $\mathcal{F} = \{\emptyset, \{1,2\}, \{3,4\}, \Omega\}$.
$X(\omega) = \omega$ ve $Y(\omega) = \mathbf{1}_{\{1,2\}}(\omega)$ fonksiyonlarından
hangisi rasgele değişkendir? Gerekçelendirin.
\end{alistirma}
\alcozum{%
$Y$: $\{Y \leq 0\} = \{3,4\} \in \mathcal{F}$; $\{Y \leq 1\} = \Omega \in \mathcal{F}$;
yani tüm alt kümeler $\mathcal{F}$'de. $Y$ rasgele değişkendir.\\
$X$: $\{X \leq 1\} = \{1\} \notin \mathcal{F}$. $X$ rasgele değişken değildir.}

\begin{alistirma}[Al.3.2]{A}{KDF'den Olasılık Hesabı}
$X$ rasgele değişkenin KDF'si:
\[
  F(x) = \begin{cases} 0 & x < 0 \\ x/2 & 0 \leq x < 1 \\ 1/2 + (x-1)/4 & 1 \leq x < 3 \\ 1 & x \geq 3 \end{cases}
\]
$P(X = 0)$, $P(X = 1)$, $P(0.5 < X \leq 2)$ ve $P(X = 3)$'ü bulun.
\end{alistirma}
\alcozum{%
$P(X = 0) = F(0) - F(0^-) = 0 - 0 = 0$ ($F$ 0'da sürekli).\\
$P(X = 1) = F(1) - F(1^-) = (1/2+0) - (1/2) = 0$ (sürekli).\\
$P(0.5 < X \leq 2) = F(2) - F(0.5) = (1/2 + 1/4) - 1/4 = 3/4 - 1/4 = 1/2$.\\
$P(X = 3) = F(3) - F(3^-) = 1 - (1/2 + 2/4) = 1 - 1 = 0$.
Bu KDF tamamen sürekli; dağılım mutlak süreklidir (karma bile değil).}

\begin{alistirma}[Al.3.3]{B}{Koşullu Dağılım}
$(X,Y)$ ortak PDF'si $f(x,y) = c\,xy$, $0 < x < 1$, $0 < y < 2$.
(a) $c$'yi bulun. (b) $f_{X|Y}(x \mid y)$'i bulun.
(c) $P(X > 0.5 \mid Y = 1)$'i hesaplayın.
\end{alistirma}
\alcozum{%
(a) $\int_0^1 \int_0^2 cxy \, dy \, dx = c \int_0^1 x \, dx \int_0^2 y \, dy = c \cdot \frac{1}{2} \cdot 2 = c$. $c = 1$.\\
(b) $f_Y(y) = \int_0^1 xy \, dx = y/2$. $f_{X|Y}(x \mid y) = \frac{xy}{y/2} = 2x$, $0 < x < 1$. ($y$'den bağımsız!)\\
(c) $P(X > 0.5 \mid Y = 1) = \int_{0.5}^1 2x \, dx = [x^2]_{0.5}^1 = 1 - 0.25 = 0.75$.}

\begin{alistirma}[Al.3.4]{B}{Jacobian Dönüşümü}
$X \sim \mathrm{Uniform}(0,1)$. $Y = -\ln X$ dönüşümü ile $Y$'nin dağılımını bulun.
\end{alistirma}
\alcozum{%
$g(x) = -\ln x$ azalan; ters $h(y) = e^{-y}$, $|h'(y)| = e^{-y}$.\\
$f_X(x) = 1$ ($0 < x < 1$), $x$ yerine $e^{-y}$ yazılırsa:
$f_Y(y) = f_X(e^{-y}) \cdot e^{-y} = 1 \cdot e^{-y} = e^{-y}$, $y > 0$.\\
$Y \sim \mathrm{Exp}(1)$. Bu quantile dönüşümü yönteminin temelidir:
herhangi bir dağılımdan örnek üretmek için KDF'nin tersini kullanmak.}

% ---

\alistirmalar

\begin{alistirma}[Al.3.5]{A}{KDF Özellikleri}
$F(x) = \frac{1}{1+e^{-x}}$ (lojistik fonksiyon) bir KDF midir?
Üç temel özelliği kontrol edin. Bu KDF mutlak sürekli bir dağılımı mı temsil eder?
PDF'yi bulun.
\end{alistirma}
\alcozum{%
Monoton artan: $F'(x) = e^{-x}/(1+e^{-x})^2 > 0$. $\checkmark$\\
Sağdan sürekli: $F$ her yerde türevlenebilir, dolayısıyla sürekli. $\checkmark$\\
Sınırlar: $F(-\infty) = 0$, $F(+\infty) = 1$. $\checkmark$\\
PDF: $f(x) = F'(x) = e^{-x}/(1+e^{-x})^2$. Bu lojistik dağılımdır.}

\begin{alistirma}[Al.3.6]{A}{PMF Doğrulama}
$p(k) = c \cdot (1/2)^k$, $k = 1, 2, 3, \ldots$
(a) $c$'yi bulun. (b) $P(X > 3)$'ü hesaplayın. (c) $P(X \text{ çift})$'i bulun.
\end{alistirma}
\alcozum{%
(a) $\sum_{k=1}^\infty c(1/2)^k = c \cdot (1/2)/(1-1/2) = c = 1$.\\
(b) $P(X > 3) = \sum_{k=4}^\infty (1/2)^k = (1/2)^4/(1-1/2) = 1/8$.\\
(c) $P(\text{çift}) = \sum_{k=1}^\infty (1/2)^{2k} = (1/4)/(1-1/4) = 1/3$.}

\begin{alistirma}[Al.3.7]{A}{Marjinal Dağılımlar}
$(X, Y)$ ortak PMF: $p(0,0)=0.1$, $p(0,1)=0.2$, $p(1,0)=0.3$, $p(1,1)=0.4$.
Marjinal PMF'leri $p_X$ ve $p_Y$'yi bulun. $X$ ve $Y$ bağımsız mıdır?
\end{alistirma}
\alcozum{%
$p_X(0) = 0.3$, $p_X(1) = 0.7$. $p_Y(0) = 0.4$, $p_Y(1) = 0.6$.\\
Bağımsız mı? $p_X(0) \cdot p_Y(0) = 0.3 \cdot 0.4 = 0.12 \neq 0.1 = p(0,0)$. Hayır.}

\begin{alistirma}[Al.3.8]{B}{Evrişim}
$X, Y$ bağımsız, $X \sim \mathrm{Uniform}(0,1)$, $Y \sim \mathrm{Uniform}(0,1)$.
$Z = X + Y$'nin PDF'sini evrişim yöntemiyle bulun ve grafiğini çizin.
\end{alistirma}
\alcozum{%
$f_Z(z) = \int_0^1 f_X(z-y) f_Y(y) \, dy = \int_0^1 \mathbf{1}_{[0,1]}(z-y) \, dy$.\\
$0 \leq z \leq 1$: $z-y \in [0,1]$ için $y \in [z-1, z] \cap [0,1] = [0, z]$; $f_Z(z) = z$.\\
$1 < z \leq 2$: $y \in [z-1, 1]$; $f_Z(z) = 2-z$.\\
Üçgen dağılım: $f_Z(z) = z$ (artan), $2-z$ (azalan). Merkezi Limit Teoremi'nin habercisi.}

\begin{alistirma}[Al.3.9]{B}{Dönüşüm: Karekök}
$X \sim \mathrm{Exp}(\lambda)$. $Y = \sqrt{X}$'in PDF'sini bulun.
$E[Y]$'yi $f_Y$ üzerinden ifade edin (Bölüm~\ref{chp:beklenti}'de hesaplanacak).
\end{alistirma}
\alcozum{%
$g(x) = \sqrt{x}$ artan; $h(y) = y^2$, $|h'(y)| = 2y$.\\
$f_X(x) = \lambda e^{-\lambda x}$ ($x > 0$).\\
$f_Y(y) = \lambda e^{-\lambda y^2} \cdot 2y = 2\lambda y e^{-\lambda y^2}$, $y > 0$.\\
Bu Rayleigh dağılımıdır (parametre $1/\sqrt{2\lambda}$).}

\begin{alistirma}[Al.3.10]{B}{Bağımsızlık Kontrolü}
$(X,Y)$ PDF'si $f(x,y) = e^{-x-y}$, $x > 0$, $y > 0$. $X$ ve $Y$ bağımsız mıdır?
Aynı soruyu $f(x,y) = 2e^{-x-y}$, $0 < y < x < \infty$ için yanıtlayın.
\end{alistirma}
\alcozum{%
Birinci: $f(x,y) = e^{-x} \cdot e^{-y} = f_X(x) \cdot f_Y(y)$. Evet, bağımsız.\\
İkinci: $f_X(x) = \int_0^x 2e^{-x-y} dy = 2e^{-x}(1-e^{-x})$. $f_Y(y) = \int_y^\infty 2e^{-x-y} dx = 2e^{-2y}$.
$f_X(x) \cdot f_Y(y) \neq f(x,y)$. Hayır, bağımsız değil.}

\begin{alistirma}[Al.3.11]{B}{Çok Boyutlu Jacobian}
$(X, Y)$ bağımsız $\Normal(0,1)$. $U = X + Y$, $V = X - Y$ dönüşümü ile
$(U, V)$'nin ortak dağılımını bulun. $U$ ve $V$ bağımsız mıdır?
\end{alistirma}
\alcozum{%
Ters: $x = (u+v)/2$, $y = (u-v)/2$. $J = \partial(x,y)/\partial(u,v)$:
$\partial x/\partial u = 1/2$, $\partial x/\partial v = 1/2$,
$\partial y/\partial u = 1/2$, $\partial y/\partial v = -1/2$.
$|\det J| = |{-1/4-1/4}| = 1/2$.\\
$f_{U,V}(u,v) = f_X((u+v)/2) f_Y((u-v)/2) \cdot 1/2
= \frac{1}{2\pi} e^{-u^2/4} e^{-v^2/4} \cdot \frac{1}{2}
= \frac{1}{4\pi} e^{-(u^2+v^2)/4}$.\\
$U \sim \Normal(0, 2)$, $V \sim \Normal(0, 2)$ ve çarpımlanıyor: bağımsızlar.}

\begin{alistirma}[Al.3.12]{C}{İspat: Bağımsız Değişkenlerin Toplamının KDF'si}
$X, Y$ bağımsız rasgele değişkenler, $Z = X + Y$ olsun.
(a) $F_Z(z) = \int_{-\infty}^{+\infty} F_X(z - y) \, dF_Y(y)$ olduğunu ispatlayın
(Lebesgue-Stieltjes integrali).
(b) Mutlak sürekli durumda bu formülün evrişim formülüne nasıl indirgeneceğini gösterin.
\end{alistirma}
\alcozum{%
(a) Bağımsızlıktan:
$F_Z(z) = P(X+Y \leq z) = E[\mathbf{1}_{X+Y \leq z}]
= E_Y[E_X[\mathbf{1}_{X \leq z-Y}]] = \int P(X \leq z-y) \, dF_Y(y) = \int F_X(z-y) \, dF_Y(y)$.\\
(b) Sürekli $Y$: $dF_Y(y) = f_Y(y) \, dy$. O zaman $F_Z(z) = \int F_X(z-y) f_Y(y) \, dy$.
$z$'ye göre türev (Leibniz): $f_Z(z) = \int f_X(z-y) f_Y(y) \, dy = (f_X * f_Y)(z)$.}

\begin{alistirma}[Al.3.13]{B}{Olasılık İntegral Dönüşümü}
$X \sim \Gama(\alpha, \beta)$. $U = F_X(X)$ dönüşümünü tanımlayın.
(a) $U$'nun dağılımını bulun.
(b) Geri dönüşüm: $U \sim \Uniform(0,1)$ verildiğinde $\Gama(2,1)$ dağılımı
nasıl simüle edersiniz? (İpucu: Üstel dağılım üzerinden)
(c) OID'nin QQ-plot oluşturmadaki rolünü açıklayın.
\end{alistirma}
\alcozum{%
(a) OID (Teorem~\ref{thm:oit}): $F_X(X) \sim \Uniform(0,1)$. Sürekli dağılım için geçerli; $\Gama$ sürekli. $U \sim \Uniform(0,1)$.\\
(b) $\Gama(2,1) = \Exp(1)+\Exp(1)$ (iki bağımsız Üstel toplamı).
$U_1, U_2 \sim \Uniform(0,1)$ bağımsız.
$X = -\ln(1-U_1) - \ln(1-U_2) = -\ln((1-U_1)(1-U_2))$.
Bu $\Gama(2,1)$ dağılımını verir.\\
(c) Veri $\{X_i\}$ için $\{F_0^{-1}(i/(n+1))\}$ teorik kuantiller ile sıralı örneklem karşılaştırılır.
$F_0 = F_X$ ise noktalar doğrusal dizilir.}

\begin{alistirma}[Al.3.14]{B}{İki Boyutlu Normal}
$(X_1, X_2) \sim \Normal_2\!\left((0,0)^\top, \begin{pmatrix}1 & \rho\\\rho & 1\end{pmatrix}\right)$.
(a) $X_1 + X_2$ ve $X_1 - X_2$'nin dağılımını bulun.
(b) $X_1 + X_2$ ve $X_1 - X_2$ bağımsız mıdır?
(c) $X_1 \mid X_2 = 1$'in dağılımını verin.
\end{alistirma}
\alcozum{%
(a) Lineer dönüşüm: $A = \begin{pmatrix}1&1\\1&-1\end{pmatrix}$, $\boldsymbol{\Sigma} = \begin{pmatrix}1&\rho\\\rho&1\end{pmatrix}$.
$A\boldsymbol{\Sigma}A^\top = \begin{pmatrix}2+2\rho & 0 \\ 0 & 2-2\rho\end{pmatrix}$.
$X_1+X_2 \sim \Normal(0, 2+2\rho)$; $X_1-X_2 \sim \Normal(0, 2-2\rho)$.\\
(b) Kovaryans sıfır (matris köşegen) $\Rightarrow$ Normal için bağımsız. Evet, bağımsız.\\
(c) $X_1|X_2=1 \sim \Normal(\rho, 1-\rho^2)$.}

\begin{alistirma}[Al.3.15]{C}{Kuantil Fonksiyon — Karma Dağılım}
$X$ karma dağılım: $P(X = 0) = 1/2$, $P(X > 0) = 1/2$ ve $X \mid X>0 \sim \Exp(1)$.
(a) $F_X(x)$'i her $x$ için yazın.
(b) $F_X^{-1}(u)$'yu $u \in (0,1)$ için hesaplayın.
(c) $U \sim \Uniform(0,1)$ kullanarak $X$ dağılımından nasıl örneklem alırsınız?
\end{alistirma}
\alcozum{%
(a) $x < 0$: $F(x) = 0$. $x = 0$: $F(0) = 1/2$. $x > 0$: $F(x) = 1/2 + (1/2)(1-e^{-x})$.\\
(b) $u \leq 1/2$: $F^{-1}(u) = 0$. $u > 1/2$: $F(x) = u \Rightarrow 1-e^{-x} = 2u-1 \Rightarrow x = -\ln(2-2u)$.
$F^{-1}(u) = -\ln(2(1-u))$ ($u > 1/2$).\\
(c) $U \sim \Uniform(0,1)$: $U \leq 1/2$ ise $X = 0$; $U > 1/2$ ise $X = -\ln(2(1-U))$.}

% =====================================================================
\endinput
